% ============================================================
% sec/06etica.tex
%
% DEBE CONTENER: los riesgos identificados y CÓMO se piensan
% mitigar. El momento de detectarlos es ANTES de entrenar.
%
% En la Etapa 3 se retoma el tema, con foco en el despliegue
% del agente.
% ============================================================
\section{Consideraciones Éticas Preliminares}
\label{sec:etica}

\subsection{Sensibilidad de los datos}
\label{subsec:sensibilidad}
\pc{¿Contienen datos personales, sensibles o protegidos? ¿Bajo qué
licencia se publicaron y qué permite esa licencia? Si no hay datos
personales, decirlo explícitamente en vez de omitir el apartado.}

\subsection{Riesgo de sesgo}
\label{subsec:sesgo}
\pc{Qué poblaciones, condiciones o regímenes están sub-representados
en la muestra, y qué falla podría causar eso aguas abajo. Declarar si
la muestra es o no aleatoria respecto del universo que el problema
sugiere.}

\subsection{Límites de generalización declarados}
\label{subsec:limitesgeneralizacion}
\pc{A qué población NO debe extrapolarse el modelo. Este apartado
existe para que nadie use el resultado fuera del alcance en el que fue
validado.}

\subsection{Mitigación prevista antes de entrenar}
\label{subsec:mitigacionetica}
\begin{itemize}
  \item \pc{Mitigación 1}
  \item \pc{Mitigación 2}
  \item \pc{Mitigación 3}
\end{itemize}

\subsection{Diferido a la Etapa 3}
\label{subsec:diferido}
La ética del despliegue, los escenarios de uso indebido y los riesgos
operativos del agente inteligente se abordan en la etapa final.